% sections/04_preprocessing_pipeline.tex  –  Preprocessing Pipeline
% \input{} into main.tex; do NOT wrap in \documentclass or \begin{document}

\section{Preprocessing Pipeline}
\label{sec:preproc}

Before any learning algorithm receives data, every $S_{11}$ trace must pass
through a deterministic preprocessing pipeline that ensures correctness,
consistency, and traceability.  This section describes \emph{Stage~1} of a
three-stage data-conditioning workflow; \Cref{sec:norm} and
\Cref{sec:augment} cover the subsequent stages.

%% ------------------------------------------------------------------
\subsection{Three-Stage Pipeline Overview}
\label{sec:preproc:stages}

The full data-conditioning workflow is decomposed into three strictly
sequential stages, each with a distinct purpose:

\begin{enumerate}[label=\textbf{Stage~\arabic*:},leftmargin=3.5em]
  \item \textbf{Validity preprocessing} (this section).
        Parses raw files, verifies structural and numerical integrity,
        quarantines malformed or outlier traces, and assigns canonical
        sample identifiers.  At the end of this stage every surviving
        sample is a well-formed numerical array with a unique ID and a
        verified class label.

  \item \textbf{Normalization} (\Cref{sec:norm}).
        Transforms feature values (magnitude, phase, or pixel
        intensities) to distributions that accelerate and stabilise
        learning.  Normalization statistics are computed on the training
        partition \emph{only} and applied identically to validation and
        test partitions.

  \item \textbf{Augmentation} (\Cref{sec:augment}).
        Synthetically expands the training set with physics-aware
        perturbations at training time only.  Augmentation is never
        applied to validation or test data.
\end{enumerate}

\noindent
\Cref{fig:pipeline_stages} sketches the three stages and their
data-flow boundaries.

\begin{figure}[ht]
\centering
\begin{tikzpicture}[
    node distance=1.2cm and 0.6cm,
    block/.style={rectangle, draw, rounded corners,
                  minimum width=3.8cm, minimum height=1.0cm,
                  align=center, font=\small},
    arrow/.style={-Stealth, thick}
]
  \node[block, fill=blue!10]   (raw)  {Raw TSV files\\(\num{1603} traces)};
  \node[block, fill=green!10,  right=of raw]  (s1)
        {Stage 1\\Validity preprocessing};
  \node[block, fill=orange!10, right=of s1]   (s2)
        {Stage 2\\Normalization};
  \node[block, fill=red!10,    right=of s2]   (s3)
        {Stage 3\\Augmentation\\(train only)};
  \draw[arrow] (raw) -- (s1);
  \draw[arrow] (s1)  -- (s2);
  \draw[arrow] (s2)  -- (s3);
\end{tikzpicture}
\caption{Three-stage data-conditioning pipeline.  Stages~2 and~3 are
         covered in \Cref{sec:norm,sec:augment}, respectively.}
\label{fig:pipeline_stages}
\end{figure}

The remainder of this section details every sub-step of Stage~1.

%% ------------------------------------------------------------------
\subsection{Touchstone File Parsing}
\label{sec:preproc:parsing}

Each $S_{11}$ trace is stored as a plain-text file with three
tab-separated columns and no header row:

\begin{center}
\begin{tabular}{ccc}
\toprule
\textbf{Column 1} & \textbf{Column 2} & \textbf{Column 3} \\
\midrule
Frequency (Hz)     & $|S_{11}|$ (dB)    & $\angle S_{11}$ (deg) \\
\bottomrule
\end{tabular}
\end{center}

\paragraph{Parsing assumptions.}
\begin{itemize}
  \item \textbf{Delimiter:} horizontal tab (\verb|\t|).
        The parser must not silently accept comma- or space-delimited
        variants; mismatched delimiters indicate a file-format error.
  \item \textbf{No header row:} the first line is expected to be numeric
        data.  However, a small number of files may carry an optional
        Touchstone-style option line beginning with
        \texttt{\#}~\cite{touchstone_spec}.  The parser shall skip any
        line whose first non-whitespace character is~\texttt{\#} and log
        the event.
  \item \textbf{Column count:} exactly three numeric fields per
        non-comment line.  Lines with fewer or more fields trigger a
        malformed-row error.
  \item \textbf{Encoding:} UTF-8 or ASCII; reject files with binary or
        non-decodable content.
\end{itemize}

\paragraph{Column interpretation.}
After parsing, each trace is a matrix
$\mathbf{T}\in\mathbb{R}^{4001\times 3}$ with columns
$(f,\;|S_{11}|_{\text{dB}},\;\angle S_{11})$.
The frequency column is in hertz; for readability in plots and tables we
convert to gigahertz (\SI{}{\giga\hertz}) by dividing by $10^{9}$,
yielding a range of approximately \SIrange{8.0}{12.0}{\giga\hertz} with
a step of roughly \SI{1}{\mega\hertz}.
The magnitude column is already in decibels (a logarithmic scale) and
requires no further transformation at this stage.
The phase column is in degrees and is retained as-is; conversion to
radians is deferred to the normalization stage if required by a specific
algorithm.

%% ------------------------------------------------------------------
\subsection{Frequency-Grid Validation}
\label{sec:preproc:freqgrid}

A fundamental assumption of the downstream pipeline is that every trace
shares the \emph{same} $N=4001$-point frequency vector.  This enables
treating each frequency bin as a fixed feature across all samples.  The
validation step proceeds as follows:

\begin{enumerate}
  \item Parse the frequency column of the first valid file and designate
        it the \emph{reference grid}~$\mathbf{f}_{\text{ref}}$.
  \item For every subsequent file, compute the element-wise absolute
        difference
        $\Delta_i = |f_i - f_{\text{ref},i}|$ for
        $i=1,\dots,4001$.
  \item If $\max_i \Delta_i > \epsilon$ (with
        $\epsilon = \SI{1}{\hertz}$), the file is flagged as
        \emph{grid-mismatched} and quarantined with a diagnostic log
        entry.
  \item Files with $N \neq 4001$ points are similarly quarantined.
\end{enumerate}

Grid mismatches could indicate a different VNA frequency plan, a
truncated sweep, or a data-export error.  Quarantined files are excluded
from all subsequent processing but retained on disk for manual
inspection.

%% ------------------------------------------------------------------
\subsection{Missing-Value and Malformed-File Handling}
\label{sec:preproc:missing}

\paragraph{Missing and non-finite values.}
Each parsed numeric entry is checked for \texttt{NaN}, $\pm\infty$, and
non-numeric content (e.g., text strings in a numeric column).
The handling policy is:

\begin{itemize}
  \item \textbf{Isolated NaN/Inf} (fewer than $0.1\,\%$ of points in a
        trace): replace via linear interpolation from the two nearest
        valid neighbours in the frequency dimension.  Log the location
        and original value.
  \item \textbf{Widespread NaN/Inf} ($\geq 0.1\,\%$ of points): reject
        the entire file.  Interpolating large gaps risks introducing
        artefacts that distort class-discriminative spectral features.
\end{itemize}

\paragraph{Malformed files.}
The following conditions result in immediate file rejection:

\begin{itemize}
  \item Fewer than $4001$ non-comment lines (truncated acquisition).
  \item Any non-comment line with a column count $\neq 3$.
  \item File size of zero bytes or failure to decode as text.
\end{itemize}

\noindent
Every rejection is logged with the file path, the reason, and a
timestamp.  The total number of rejected files is reported before
proceeding.

%% ------------------------------------------------------------------
\subsection{Coordinate Transformation (3-Channel)}
\label{sec:preproc:transform}

After parsing the native dB/Phase data and verifying structural
integrity, the pipeline expands the two-column feature space into
three channels that form the primary Route~A representation
(see~\Cref{sec:repr:routeA:hybrid}):

\begin{description}
  \item[Channel~1:] $|S_{11}|_{\text{dB}}$ — passed directly from
        column~2 of the raw file.  Retains the full dynamic range of
        the magnitude (critical for the deep $\SI{-62}{\decibel}$
        resonance).
  \item[Channel~2:]
        $\text{Re} = 10^{\,|S_{11}|_{\text{dB}}\,/\,20} \cdot
        \cos\!\bigl(\angle S_{11}\cdot\pi/180\bigr)$.
  \item[Channel~3:]
        $\text{Im} = 10^{\,|S_{11}|_{\text{dB}}\,/\,20} \cdot
        \sin\!\bigl(\angle S_{11}\cdot\pi/180\bigr)$.
\end{description}

\noindent
This transformation is a deterministic, lossless, invertible mapping
that introduces no free parameters and therefore has no data-leakage
risk.  It is applied identically to all partitions (train, validation,
test).  After this step, each trace is stored as a matrix
$\mathbf{X} \in \mathbb{R}^{4001 \times 3}$.

%% ------------------------------------------------------------------
\subsection{Outlier Inspection}
\label{sec:preproc:outliers}

Outlier screening is performed \emph{before} training to detect
anomalous traces that may represent measurement errors, sensor
malfunctions, or mislabelled files.  The procedure operates
per-frequency-bin within each class:

\begin{enumerate}
  \item For each class $c \in \{\texttt{dry},\,\texttt{wet},\,
        \texttt{ice}\}$ and each frequency index $i=1,\dots,4001$,
        compute the mean $\mu_{c,i}$ and standard deviation
        $\sigma_{c,i}$ of the magnitude values across all samples in
        class~$c$.
  \item Flag any sample whose magnitude at bin~$i$ satisfies
        \[
            |x_{i} - \mu_{c,i}| > 4\,\sigma_{c,i}
        \]
        as a candidate outlier at that bin.
  \item A sample flagged at more than $1\,\%$ of frequency bins is
        promoted to a \emph{trace-level outlier candidate}.
  \item Trace-level candidates are \emph{not} automatically removed;
        they are presented for visual inspection via class-aggregated
        envelope plots (mean $\pm\;2\sigma$ and $\pm\;4\sigma$ bands),
        and the analyst decides on retention, relabelling, or exclusion.
\end{enumerate}

The same procedure is repeated for the phase channel.  Automatic
removal without human review is deliberately avoided because a trace
that appears statistically anomalous may represent a genuine edge case
that the classifier should learn.

\begin{tcolorbox}[
  colback=red!5!white,
  colframe=red!75!black,
  title={\textbf{Anti-Leakage Requirement}},
  fonttitle=\bfseries
]
When the outlier screening is used to set automatic rejection
thresholds (rather than visual inspection), the statistics
$(\mu_{c,i},\,\sigma_{c,i})$ \textbf{must be fitted exclusively on
the training fold}.  Validation and test folds are evaluated against
these frozen thresholds.  Computing $(\mu,\sigma)$ on the full
\num{1603}-sample dataset prior to splitting constitutes data leakage
and would produce optimistically biased performance estimates (see
also~\Cref{sec:norm:leakage}).
\end{tcolorbox}

%% ------------------------------------------------------------------
\subsection{Class-Label Mapping}
\label{sec:preproc:labels}

Class labels are derived from the folder structure of the dataset:

\begin{center}
\begin{tabular}{lcc}
\toprule
\textbf{Folder name} & \textbf{Integer label} & \textbf{Sample count} \\
\midrule
\texttt{dry}  & 0 & 601 \\
\texttt{wet}  & 1 & 501 \\
\texttt{ice}  & 2 & 501 \\
\midrule
\multicolumn{2}{l}{\textbf{Total}} & \textbf{1603} \\
\bottomrule
\end{tabular}
\end{center}

The class imbalance ratio is $601:501:501 \approx 1.20:1:1$, which is
mild but non-trivial; it is accounted for during evaluation
(\Cref{sec:eval}) via macro-averaged metrics.

\paragraph{Trace-numbering gaps.}
Inspection of the raw file names reveals numbering gaps in the trace
indices: traces \texttt{Trc601}--\texttt{Trc699} and
\texttt{Trc1201}--\texttt{Trc1449} are absent across all classes.
These gaps are noted as potential missing acquisitions or intentional
exclusions by the data provider.  The pipeline does not attempt to
impute or reconstruct the missing trace indices; it simply uses the
files that are present.

%% ------------------------------------------------------------------
\subsection{Deterministic Sample Indexing}
\label{sec:preproc:indexing}

To ensure full reproducibility across experiments, each surviving sample
is assigned a canonical integer identifier:

\begin{enumerate}
  \item Sort all valid file paths lexicographically within each class
        folder.
  \item Concatenate the sorted lists in label order
        ($0 \to 1 \to 2$).
  \item Assign sequential IDs starting from~$0$.
\end{enumerate}

A mapping file \texttt{sample\_index.csv} is written with columns
\texttt{sample\_id}, \texttt{class\_label}, \texttt{class\_name}, and
\texttt{original\_filepath}.  This file is version-controlled alongside
the code and referenced in all subsequent logs to ensure that any
experimental result can be traced back to the exact input files.

%% ------------------------------------------------------------------
\subsection{K-Means Separability Gate}
\label{sec:preproc:kmeans_gate}

Before committing to supervised training, the pipeline applies a
\emph{K-Means separability gate} that verifies whether the
three-channel feature representation produces naturally separable
clusters aligned with the three physical states.  This probe is
analogous to the Level~0 evaluation described in
\Cref{sec:algorithms}, but positioned earlier in the pipeline as a
quality checkpoint.

\paragraph{Procedure.}
\begin{enumerate}
  \item Run K-Means with $k = 3$ on the three-channel feature vectors,
        using $\geq 50$ random initialisations.
  \item Compute the Adjusted Rand Index (ARI) between the discovered
        clusters and the known class labels.
  \item If $\text{ARI} < 0.5$, the representation is flagged for
        revision before proceeding to supervised model training.
\end{enumerate}

\begin{tcolorbox}[
  colback=red!5!white,
  colframe=red!75!black,
  title={\textbf{Anti-Leakage Requirement}},
  fonttitle=\bfseries
]
The K-Means clustering model \textbf{must be fitted only on the
training split}.  Applying K-Means to the entire \num{1603}-sample
dataset prior to splitting constitutes \emph{representation leakage}:
the cluster centroids encode distributional information from
validation and test samples, making the separability assessment
unreliable as a predictor of generalisation performance.
\end{tcolorbox}

%% ------------------------------------------------------------------
\subsection{Stage~1 Summary}
\label{sec:preproc:summary}

\Cref{tab:preproc_checklist} summarises the validity checks applied
during Stage~1.

\begin{table}[ht]
\centering
\caption{Stage~1 validity-preprocessing checklist.}
\label{tab:preproc_checklist}
\small
\begin{tabular}{p{4.0cm} p{4.5cm} p{3.8cm}}
\toprule
\textbf{Check} & \textbf{Condition} & \textbf{Action on failure} \\
\midrule
Comment-line filtering &
  Line starts with \texttt{\#} &
  Skip line, log event \\
Column count &
  $\neq 3$ fields per line &
  Reject file \\
Row count &
  $\neq 4001$ data rows &
  Quarantine file \\
Frequency grid &
  $\max|\Delta f| > \SI{1}{\hertz}$ vs.\ reference &
  Quarantine file \\
Non-finite values &
  NaN or $\pm\infty$ entries &
  Interpolate ($<0.1\,\%$) or reject \\
Outlier screening &
  $>4\sigma$ in $>1\,\%$ of bins &
  Flag for visual review \\
\bottomrule
\end{tabular}
\end{table}

After Stage~1, the dataset consists of $N_{\text{valid}} \leq 1603$
well-formed, grid-aligned, canonically indexed traces ready for
normalization (\Cref{sec:norm}).
